\chapter{Dengue Haemorrhagic Fever}
\index{Dengue Haemorrhagic Fever}

\section*{Definition and Overview}
Dengue haemorrhagic fever (DHF) is a severe form of dengue infection, now more commonly called severe dengue, characterised by increased blood vessel permeability that allows fluid to leak from the bloodstream into body tissues. This plasma leakage causes low blood pressure, fluid accumulation in the chest and abdomen, and bleeding tendencies, and can progress to dengue shock syndrome, which is life-threatening. DHF occurs most often in people experiencing a second dengue infection with a different virus serotype. With early recognition and hospital care, the vast majority of people survive.

\section*{Epidemiology}
Severe dengue, including dengue haemorrhagic fever, occurs in a small proportion of dengue infections, estimated at around 5\% or less of symptomatic cases. The risk is highest in children in endemic countries and in people infected a second time with a different serotype. In many countries, severe dengue is rare but occurs in returned travellers and occasionally in local outbreaks. Globally, severe dengue causes tens of thousands of deaths each year, mainly in Asia and Latin America, though mortality is below 1\% with good hospital care. The geographic range and case numbers continue to rise with mosquito spread and climate change.

\section*{Aetiology and Pathophysiology}
\begin{itemize}
    \item \textbf{Second infection:} antibody-dependent enhancement, in which antibodies from a first dengue infection help a different serotype enter cells more efficiently, drives most severe disease
    \item \textbf{Plasma leakage:} the immune response damages the lining of small blood vessels, allowing fluid and proteins to leak into tissues
    \item \textbf{Haemoconcentration:} as plasma is lost, the blood becomes thicker, shown by a rising haematocrit
    \item \textbf{Low platelets:} platelet counts fall, increasing bleeding risk
    \item \textbf{Shock:} severe plasma leakage reduces circulating blood volume, causing hypotension and organ hypoperfusion
    \item \textbf{Critical phase:} leakage typically occurs as the fever subsides, around days 3--7 of illness
\end{itemize}

\section*{Clinical Features}
\begin{itemize}
    \item Initial dengue fever followed by a critical phase with warning signs
    \item Severe abdominal pain, persistent vomiting, and rapid breathing
    \item Bleeding from the gums, nose, or injection sites; blood in vomit or stool
    \item Lethargy, restlessness, or irritability
    \item Cold, clammy skin, weak rapid pulse, and low blood pressure in shock
    \item Reduced urine output
    \item Fluid accumulation in the chest (pleural effusion) or abdomen (ascites)
    \item A rising haematocrit with a falling platelet count on blood tests
\end{itemize}

\section*{Differential Diagnosis}
\begin{itemize}
    \item Other causes of fever with bleeding, including leptospirosis and other viral haemorrhagic fevers
    \item Malaria, which can cause severe illness with low platelets
    \item Meningococcal disease, a bacterial infection with rash and shock
    \item Sepsis from other causes
    \item Typhoid fever with intestinal bleeding
    \item Immune thrombocytopenia and other haematological disorders
\end{itemize}

\section*{When to Seek Medical Care}
Any person with suspected dengue who develops warning signs must be taken to hospital immediately. Warning signs include severe abdominal pain, persistent vomiting, bleeding, drowsiness, restlessness, and difficulty breathing. People at increased risk, including children, pregnant women, older adults, and those with chronic disease, should be monitored particularly closely.

\textbf{Contact your local emergency services for:}
\begin{itemize}
    \item Any signs of shock: cold skin, rapid weak pulse, drowsiness, or fainting
    \item Significant bleeding, including vomiting blood or black stools
    \item Difficulty breathing
    \item Seizures or loss of consciousness
    \item Inability to drink or keep down any fluids
\end{itemize}

\section*{Diagnosis and Investigations}
\begin{itemize}
    \item \textbf{Clinical criteria:} fever plus warning signs and evidence of plasma leakage, bleeding, or shock
    \item \textbf{Full blood count:} falling platelets and rising haematocrit are the key markers of plasma leakage
    \item \textbf{Ultrasound:} detects fluid in the chest or abdomen
    \item \textbf{Blood tests:} liver function, coagulation studies, and electrolytes
    \item \textbf{Virological testing:} PCR, NS1 antigen, and serology confirm dengue
    \item \textbf{Close monitoring:} blood pressure, pulse, urine output, and haematocrit are monitored frequently in hospital
\end{itemize}

\section*{Management}
\begin{itemize}
    \item \textbf{Hospital admission:} all suspected cases of severe dengue require hospital care
    \item \textbf{Intravenous fluids:} careful, controlled fluid replacement to maintain circulation without overloading the heart and lungs
    \item \textbf{Fluid balance monitoring:} strict recording of intake, output, and vital signs, with regular blood tests
    \item \textbf{Blood transfusions:} for significant bleeding; platelets may be given for severe thrombocytopenia with bleeding
    \item \textbf{Shock management:} intensive care with fluids and medications to support blood pressure
    \item \textbf{Monitoring for fluid overload:} intravenous fluids are reduced once plasma leakage stops, which is signalled by the haematocrit falling
    \item \textbf{Supportive care:} fever control with paracetamol, rest, and careful observation
\end{itemize}

\section*{Complications}
\begin{itemize}
    \item Dengue shock syndrome, a life-threatening fall in blood pressure
    \item Severe bleeding, including gastrointestinal bleeding
    \item Fluid overload and pulmonary oedema from over-rapid fluid replacement
    \item Organ failure affecting the liver, kidneys, or brain
    \item Encephalitis or encephalopathy
    \item Death in untreated or delayed cases
    \item Prolonged recovery with fatigue in survivors
\end{itemize}

\section*{Prognosis and Outcomes}
With early hospital care and careful fluid management, the case fatality rate of severe dengue is under 1\%, compared with 10--20\% or higher without treatment. The critical phase typically lasts 24--48 hours, after which the leaking vessels recover and the person improves rapidly. Most survivors make a full recovery, although fatigue may persist for weeks. The outcome depends on early recognition of warning signs, prompt admission, and expert fluid management, which is why public education about warning signs is a cornerstone of dengue control.

\section*{Prevention and Self-Care}
\begin{itemize}
    \item Prevent dengue infection by avoiding mosquito bites in affected areas
    \item In dengue-endemic regions, seek medical care early for any fever
    \item Know the warning signs of severe dengue and act on them immediately
    \item Do not use aspirin or ibuprofen if dengue is suspected
    \item Encourage children and high-risk people to be assessed early
    \item Support mosquito control in the community by removing breeding sites
    \item Consider vaccination only on specialist advice, as it is not suitable for everyone
\end{itemize}

\section*{Sources and Further Reading}
The following reputable sources provide further information for healthcare students and patients seeking reliable, current advice about dengue haemorrhagic fever.

\begin{itemize}
    \item World Health Organization. \textit{Dengue and severe dengue: clinical management guidelines.} https://www.who.int/
    \item Centers for Disease Control and Prevention. \textit{Severe dengue: clinical guidance.} https://www.cdc.gov/dengue/
    \item Queensland Health. \textit{Dengue haemorrhagic fever information.} https://www.health.qld.gov.au/
    \item MSD Manual. \textit{Dengue: professional overview and management.} https://www.msdmanuals.com/
    \item MedlinePlus. \textit{Dengue haemorrhagic fever.} https://medlineplus.gov/
    \item NSW Health. \textit{Dengue: public health information.} https://www.health.nsw.gov.au/
\end{itemize}
